% Standalone appendix document - compiles on its own:
%   latexmk -pdf -outdir=build 01-development-guidelines.tex
\documentclass[a4paper,11pt]{article}
\usepackage[english]{babel}
\usepackage{../../appendix-style}
\usepackage{listings}
%\usepackage{minted}

\lstset{basicstyle=\ttfamily\small, columns=fullflexible, breaklines=true}

\title{10.1 Development Guidelines}
\author{SW4PRJ4 - Group 3}
\date{\today}

\begin{document}
\maketitle

\begin{versionhistory}
  1.0 & 09-09-2026 & Group 3 & LaTeX templates \\
  1.1 & 16-09-2026 & Group 3 & Added local setup, git workflow and group rules \\
  1.2 & 18-09-2026 & Group 3 & Added reference to appendix 10.3 \\
  1.3 & 23-09-2026 & Group 3 & Added sprints and the Kaneo board \\
  1.4 & 28-10-2026 & Group 3 & Simpler text, updated AI rules and epic labels \\
\end{versionhistory}

This appendix collects the practical rules the group works by. The reasons behind the technology choices are in appendix 2.1.

\section{LaTeX Templates and Styles}
We have made templates for the documents, so they look the same and a new document is quick to start. In an appendix you only need these two lines:
\begin{lstlisting}[language={[LaTeX]TeX}]
\usepackage[english]{babel} %english or danish to determine language
\usepackage{../../appendix-style} % to use appendix template
\end{lstlisting}

The report and appendices are written in English. Agendas and meeting minutes are written in Danish.

\section{Local Setup}
You need the .NET 10 SDK, Node.js and Docker.

\begin{lstlisting}
dotnet build                 # build the backend
dotnet test                  # run unit and integration tests
docker compose up            # start api, db and keycloak

cd clients
npm install                  # install all client packages (also activates git hooks)
npm test                     # Vitest in all client packages
npm run typecheck            # tsc in all client packages
npm run generate:api-client  # regenerate api-client from the OpenAPI document
\end{lstlisting}

When compose is running, the API is at localhost:5001 and Keycloak at localhost:8082.

\section{Configuration and Secrets}
Passwords and other secrets are never committed. They are set as environment variables or in a local .env file, which git ignores. The file .env.example shows which variables exist. Locally you need at least the database password and the Keycloak admin password.

In production the variables are set in Coolify. The production compose file stops with an error if one is missing, so we do not deploy with an empty password by accident.

\section{Git Workflow}
We use GitHub Flow. Why we chose it is described in appendix 2.1.

\subsection{Branches}
The main branch is the only long-lived branch, and it must always work. All work happens on a short-lived branch created from an up-to-date main branch:
\begin{lstlisting}
git switch main
git pull
git switch -c feat/pmu-12-household-profile
\end{lstlisting}

A branch name is a type and a slug with a slash between them, like in the example. The slug is lowercase English words separated by hyphens, and starts with the Kaneo task number when there is one. A branch should live for days, not weeks.

\subsection{Types}
The same types are used for branches and commits:

\begin{center}
\begin{tabular}{@{}l l@{}}
\textbf{Type} & \textbf{Used for} \\
\hline
\texttt{feat} & a new feature \\
\texttt{fix} & a bug fix \\
\texttt{refactor} & restructuring code without changing behaviour \\
\texttt{test} & only tests \\
\texttt{docs} & report, appendices and other documentation \\
\texttt{style} & formatting only \\
\texttt{perf} & performance improvement \\
\texttt{build} & build setup and dependencies \\
\texttt{ci} & GitHub Actions \\
\texttt{chore} & other maintenance \\
\texttt{revert} & undoing an earlier commit \\
\end{tabular}
\end{center}

\subsection{Commit Messages}
We follow Conventional Commits. The first line has a type, an optional scope in brackets and a short description. It is written in English, in the imperative (``add'', not ``added''), and it is at most 72 characters.
\begin{lstlisting}
feat(api): add household profile endpoint
fix(mobile): show error when the API is unreachable
docs: update technical analysis
\end{lstlisting}

\subsection{Pull Requests}
We keep the branch up to date by merging the main branch into it:
\begin{lstlisting}
git merge origin/main
\end{lstlisting}
We do not rebase or force push branches that others may have pulled.

When the work is ready, we open a pull request against the main branch. The title must also follow Conventional Commits, and the description links the Kaneo task. The pull request needs green CI and one review from another group member before it can be merged.

We use squash merge. The whole branch becomes one commit on the main branch, with the pull request title as the message. The branch is deleted afterwards.

\subsection{Releases}
At the end of a sprint, or before a hand-in, we put a tag on the main branch, for example v0.1.0. We do not use release branches.

\subsection{Automatic Checks}
The rules are checked automatically, so nobody has to remember them all. One git hook rejects a commit whose message does not follow the format. Another rejects a push directly to the main branch and a push of a branch with a wrong name.

The hooks are turned on when you run npm install in the clients folder. Without Node.js, run this once instead:
\begin{lstlisting}
git config core.hooksPath .githooks
\end{lstlisting}

A GitHub Actions workflow checks the branch name and the pull request title again on GitHub, where it cannot be skipped. All three use the same script, so the rules are only written in one place.

\section{Languages}
\begin{itemize}
    \item Code, names, comments and commit messages: English.
    \item All text the user sees in the apps, including error messages from the API: English (NFR-01).
    \item Kaneo tasks: Danish.
\end{itemize}

\section{AI Tools}
\begin{itemize}
    \item The shared rules for the AI tools are in the AGENTS.md file in the repository root. Change a rule there, not in a tool-specific file.
    \item The AI may read the Kaneo board, create and comment on tasks, create, add and remove labels, link tasks to each other, and set the start and due date of a task to the first and last day of its sprint. New tasks from the AI go to the backlog.
    \item When a group member asks, the AI may move tasks between the backlog and the board, and it may estimate a task and set its story point label. The group corrects the estimate if it disagrees. Which tasks are moved, and when, is decided by the group. Assigning and deleting tasks is also done by the group.
    \item The AI cannot read .env files or key files.
    \item Whoever commits code must be able to explain it, also at the exam.
\end{itemize}

\section{Sprints and the Kaneo Board}
A sprint is two weeks, from Wednesday to the Tuesday two weeks later. The exceptions are sprint 2, a one-week mini sprint so that it ends with a demo, sprint 3, which is three weeks because the autumn break falls in the middle of it, and sprint 4, which is one week. Sprint review, retrospective and planning are held on the same day, and the sprint goal is written in the minutes of that meeting.

\begin{center}
\begin{tabular}{@{}ll@{}}
Sprint & Period \\ \hline
1 & 16/9 -- 29/9 \\
2 & 30/9 -- 6/10 \\
3 & 7/10 -- 27/10 \\
4 & 28/10 -- 3/11 \\
5 & 4/11 -- 17/11 \\
6 & 18/11 -- 1/12 \\
\end{tabular}
\end{center}

The report is handed in on Friday 11/12. The days from 2/12 to 11/12 are not a sprint, but a buffer before the hand-in.

\begin{itemize}
    \item The Backlog view in Kaneo is the product backlog. The board is the sprint backlog.
    \item A task moves from the backlog to \emph{To Do} at sprint planning, and only when it is ready. That means it has a description that starts with the requirements it covers (Krav: FR-xx), a definition of done, a story point label, an epic label and an owner.
    \item Tasks in the backlog have no sprint label. A task gets the label of a sprint when it is pulled into that sprint at sprint planning.
    \item Every task in a sprint gets the label of that sprint. A task that is not finished keeps its label and gets the next one as well, so spillover is visible.
    \item Story points use the planning poker scale: 0, \textonehalf, 1, 2, 3, 5, 8, 13, 20, 40 and ?. The group estimates together.
    \item Epic labels are short names for the epics in the requirements specification: Account, Household (order contents), Subscription (subscription and orders), Notifications, Catalogue (package catalogue) and Dispatch. The Household label is from before the requirements were changed. Work that is not a requirement gets \emph{Teknisk} or \emph{Rapport}.
\end{itemize}

\section{Definition of Done}
A task is done when:
\begin{itemize}
    \item The build and all tests are green, including the client tests if the client code was changed.
    \item The change is merged to the main branch through a pull request with green CI and one review.
    \item The Kaneo task is linked in the pull request.
    \item The person who made the change can explain it.
\end{itemize}

Appendix 10.3 describes what happens after the merge: the image build, the deployment to production and how we roll back.

\end{document}
